\begin{figure}[htbp]
\centering
\begin{tikzpicture}[
  scale=0.72, transform shape, font=\small,
  %%%%%%%%% estados da máquina
  modo/.style={draw, rounded corners=2mm, fill=white, align=center,
               minimum width=34mm, minimum height=11mm, inner sep=3pt},
  %%%%%%%%% relacións
  fluxo/.style={-{Latex[length=2.4mm, width=2mm]}},
  etiq/.style={font=\scriptsize, align=center, inner sep=2pt, fill=white}
]

\node[modo] (oam)    at (0,0)      {\textbf{Modo 2}\\ Busca en \acrshort{oam}};
\node[modo] (draw)   at (72mm,0)   {\textbf{Modo 3}\\ Transferencia de píxeles};
\node[modo] (hblank) at (72mm,-36mm) {\textbf{Modo 0}\\ \textit{HBlank}};
\node[modo] (vblank) at (0,-36mm)  {\textbf{Modo 1}\\ \textit{VBlank}};

%%%%%%%%% transicións
\draw[fluxo] (oam) -- node[etiq, above]{80 \textit{dots}} (draw);

\draw[fluxo] (draw) -- node[etiq, right]{172 \textit{dots}\\ debúxase a liña \texttt{ly}} (hblank);

\draw[fluxo] (hblank) -- node[etiq, below, pos=0.55]
      {204 \textit{dots}, \texttt{ly = 144}\\ \textit{VBlank} e cadro completo} (vblank);

\draw[fluxo] (hblank.north west) -- node[etiq, pos=0.45]
      {204 \textit{dots}, \texttt{ly < 144}} (oam.south east);

\draw[fluxo] (vblank) -- node[etiq, left]{\texttt{ly = 154}\\ \texttt{ly = 0}} (oam);

\draw[fluxo] (vblank.south west) to[out=-135, in=135, looseness=6]
      node[etiq, below left, pos=0.5]{456 \textit{dots}\\ \texttt{ly += 1}} (vblank.west);

%%%%%%%%% entrada
\node[etiq, above=5mm of oam, fill=none] (ent) {\texttt{dots += ciclos da instrución}};
\draw[fluxo] (ent) -- (oam);

\end{tikzpicture}
\caption{Máquina de estados dos modos da \acrshort{ppu} tal e como se implementa.}
\label{fig:ppu-modos}
\end{figure}
